\section{Additional Details and Results on Image Generation}
\label{sec:image_generation}


\subsection{Training-free Acceleration of Latent-space Image Generation}
\label{subsec:image_latent}

For training-free inference acceleration of image generation, we use NVIDIA A100 GPUs with 80~GB of VRAM to match RALU~\citep{ralu}'s evaluation protocol. For training-free inference acceleration of video generation, as well as image editing and ablation studies, we use NVIDIA H100 GPUs with 80~GB of VRAM. For training-free acceleration of latent-space image generation, we use FLUX.1-dev~\citep{black2024flux} as the base model and follow the evaluation protocol of RALU~\citep{ralu}. Table~\ref{tab:flux-extended} extends Table~\ref{tab:flux-main} with additional training-free acceleration baselines from RALU. We organize the comparison into three wall-clock speedup tiers:

\begin{itemize}
    \item \textbf{Small speedup}: the 50-step native full-resolution baseline, compared against RALU (50 steps) and our method (40 steps), yielding approximately $1.7\times$ speedup;
    \item \textbf{Medium speedup}: the 10-step native full-resolution baseline, used as an iso-compute reference for RALU (15 steps) and our method (12 steps), yielding approximately $5\times$ speedup; and
    \item \textbf{Large speedup}: the 7-step native full-resolution baseline, used as an iso-compute reference for RALU (10 steps) and our method (10 steps), yielding approximately $7\times$ speedup.
\end{itemize}
All speedups are calculated relative to the wall-clock latency of the 50-step FLUX.1-dev baseline on a single A100 GPU, following RALU~\citep{ralu}. As in Table~\ref{tab:flux-main}, we evaluate image quality using CLIP-IQA~\citep{wang2023clipiqa} and NIQE~\citep{mittal2013niqe}, prompt alignment using T2I-CompBench~\citep{t2icompbench} and GenEval~\citep{geneval}, and overall quality using ImageReward~\citep{imagereward}.
Following RALU's evaluation protocol, ImageReward, CLIP-IQA, and NIQE are averaged over 5,000 images generated from MS-COCO validation prompts. For compositional alignment, GenEval is evaluated with four random seeds per prompt on its object-focused prompt suite, and T2I-CompBench is evaluated on its spatial, non-spatial, and complex prompt subsets with four random seeds per prompt. Unless otherwise specified, all metrics are computed using their official or default evaluation settings.

We keep the error threshold fixed at $\delta=0.01$, use $T_\Phi = \text{DCT}$, and compare $S=2$ and $S=3$ resolution-stage settings. To align with the multi-resolution training distribution of FLUX.1-dev, we use latent dimensions $64^2$ ($s_1 = 0.5$) and $128^2$ $(s_2 = 1.0)$ (pixel dimensions $512^2$ and $1024^2$) for $S=2$, and latent dimensions $32^2$ $(s_1 = 0.25)$, $64^2$ $(s_2 = 0.5)$, and $128^2$ ($s_3 = 1.0$) (pixel dimensions $256^2$, $512^2$, and $1024^2$) for $S=3$.

By combining spectral transformations for progressive-resolution generation with our derived resolution schedule, our method achieves a better speed-quality tradeoff than the reduced-step FLUX.1-dev baselines and existing temporal and spatial diffusion acceleration methods, including progressive-resolution methods such as RALU~\citep{ralu}. In the largest speedup tier, our $S=3$ setting reaches $7.09\times$ wall-clock speedup and $7.36\times$ FLOPs speedup while preserving competitive image quality. We further provide extended qualitative baseline comparisons in Figures~\ref{fig:flux-1}--\ref{fig:flux-3}. FLUX.1-dev with reduced steps degrades image quality and exhibits over-saturation artifacts, while RALU introduces noticeable noise artifacts in reduced-step settings.


\begin{table*}[t!]
    \centering
    \caption{\textbf{Extended quantitative comparisons on training-free latent-space image generation.} We evaluate FLUX.1-dev at $1024^2$ resolution and group methods into the same $1\times$, $5\times$, and $7\times$ wall-clock speedup tiers as Table~\ref{tab:flux-main}. Rows shared with Table~\ref{tab:flux-main} use the same values. Speedup is calculated relative to the 50-step FLUX.1-dev baseline wall-clock runtime on a single A100 GPU. \textcolor{blue}{\textbf{T}} and \textcolor{red}{\textbf{S}} denote temporal and spatial acceleration methods, respectively. Baseline metrics are copied from RALU~\citep{ralu} under the same evaluation protocol.}
    \label{tab:flux-extended}
    \resizebox{\textwidth}{!}{
    \begin{scriptsize}
    \begin{tabular}{cccccccccc}
         \toprule
        \multirow{2}{*}{Method} & \multirow{2}{*}{Accel.} & \multirow{2}{*}{Latency (s) $\downarrow$} & \multirow{2}{*}{Speedup (s) $\uparrow$} & \multirow{2}{*}{TFLOPs $\downarrow$} & Overall & \multicolumn{2}{c}{Image quality} & \multicolumn{2}{c}{Text alignment} \\
        \cmidrule(lr){6-10}
        & & & & & ImageReward $\uparrow$ & CLIP-IQA $\uparrow$ & NIQE $\downarrow$ & T2I-Comp. $\uparrow$ & GenEval $\uparrow$  \\
        \midrule
        FLUX (50 steps) & - & 25.1 & 1.00$\times$ & 2991.01 & \textbf{1.095} & 0.707 & 6.75 & \textbf{0.634} & \textbf{0.698}  \\
        RALU~\citep{ralu}   & \textcolor{red}{\textbf{S}}  & 15.85 & 1.58$\times$ & \underline{1749.94} & 1.028 & \underline{0.712} & \textbf{6.07} & 0.613 & 0.648 \\
        \textbf{Ours ($S=2$)}     & \textcolor{red}{\textbf{S}}       & \underline{15.16} & \underline{1.66}$\times$ & 1755.22 & \underline{1.049} & \textbf{0.719} & 6.43 & \underline{0.617} & \underline{0.654} \\
        \textbf{Ours ($S=3$)}     & \textcolor{red}{\textbf{S}}       & \textbf{14.54} & \textbf{1.73}$\times$ & \textbf{1672.04} & 1.015 & 0.711 & \underline{6.33} & 0.593 & 0.640 \\
        \cmidrule(lr){1-10}
        FLUX (10 steps) & \textcolor{blue}{\textbf{T}} & 5.18 & 4.84$\times$ & 610.02 & 0.981 & 0.679 & 6.93 & 0.618 & 0.647 \\
        $\Delta$-DiT~\citep{deltadit}      & \textcolor{blue}{\textbf{T}} & 7.42 & 3.38$\times$ & 772.10 & 0.102 & 0.487 & 9.60 & 0.306 & 0.397 \\
        ToCa~\citep{toca}                  & \textcolor{blue}{\textbf{T}} & 15.5 & 1.62$\times$ & 601.12 & -1.827 & 0.253 & 10.6 & 0.259 & 0.137 \\
        TeaCache~\citep{teacache}          & \textcolor{blue}{\textbf{T}} & 5.23 & 4.80$\times$ & 610.59 & 0.944 & 0.665 & 7.92 & 0.620 & 0.647 \\
        TaylorSeer~\citep{taylorseer}      & \textcolor{blue}{\textbf{T}} & 9.34 & 2.69$\times$ & 556.72 & 0.972 & 0.684 & 6.77 & 0.594 & 0.619 \\
        Bottleneck~\citep{bottleneck}      & \textcolor{red}{\textbf{S}}  & 5.37 & 4.67$\times$ & 571.23 & 0.889 & 0.661 & 9.16 & 0.620 & \textbf{0.687} \\
        RALU~\citep{ralu}                  & \textcolor{red}{\textbf{S}}  & 5.04 & 4.98$\times$ & 540.47 & 1.022 & \underline{0.700} & \textbf{6.43} & \textbf{0.626} & 0.652 \\
        \textbf{Ours ($S=2$)}     & \textcolor{red}{\textbf{S}}       & \underline{4.35} & \underline{5.77}$\times$ & \underline{500.34} & \textbf{1.059} & 0.696 & 6.69 & \underline{0.624} & \underline{0.655} \\
        \textbf{Ours ($S=3$)}     & \textcolor{red}{\textbf{S}}       & \textbf{4.12} & \textbf{6.09}$\times$ & \textbf{469.15} & \underline{1.042} & \textbf{0.701} & \underline{6.53} & 0.623 & 0.637 \\
        \cmidrule(lr){1-10}
        FLUX (7 steps) & \textcolor{blue}{\textbf{T}} & 3.79 & 6.62$\times$ & 431.45 & 0.920 & 0.660 & 8.25 & 0.594 & 0.583 \\
        TeaCache~\citep{teacache}          & \textcolor{blue}{\textbf{T}} & 4.21 & 5.96$\times$ & 431.83 & 0.733 & 0.623 & 13.7 & 0.599 & 0.594 \\
        TaylorSeer~\citep{taylorseer}      & \textcolor{blue}{\textbf{T}} & 7.00 & 3.59$\times$ & 431.74 & 0.660 & 0.646 & 9.43 & 0.514 & 0.446 \\
        Bottleneck~\citep{bottleneck}      & \textcolor{red}{\textbf{S}}  & 3.78 & 6.64$\times$ & 431.52 & 0.792 & 0.631 & 8.71 & 0.605 & \underline{0.672} \\
        RALU~\citep{ralu}                  & \textcolor{red}{\textbf{S}}  & 3.75 & 6.69$\times$ & \underline{426.01} & 0.999 & 0.681 & 6.87 & \textbf{0.633} & \textbf{0.682} \\
        \textbf{Ours ($S=2$)}     & \textcolor{red}{\textbf{S}}       & \underline{3.70} & \underline{6.78}$\times$ & 427.03 & \textbf{1.039} & \underline{0.689} & \underline{6.78} & 0.620 & 0.667 \\
        \textbf{Ours ($S=3$)}     & \textcolor{red}{\textbf{S}}       & \textbf{3.54} & \textbf{7.09}$\times$ & \textbf{406.24} & \underline{1.015} & \textbf{0.694} & \textbf{5.99} & \underline{0.627} & 0.637 \\
        \bottomrule
    \end{tabular}
    \end{scriptsize}}
\end{table*}


\subsection{Spectral-transformation-based Fine-tuning for Latent- and Pixel-Space Image Generation}
\label{subsec:image_pixel}
\paragraph{Additional latent-space fine-tuning results.} Since FLUX.1-dev is a guidance-distilled model and its base model is not open-sourced, we additionally conduct latent-space image generation fine-tuning experiments on Z-Image~\citep{zimage}, which employs the same VAE architecture as FLUX.1-dev and hence the same $\delta$-optimal resolution schedules. We adopt Low-Rank Adaptation (LoRA) of rank 32 and train for 2000 steps following the spectral-transformation-based fine-tuning procedure in Sec.~\ref{sec:method:lora}. For fair comparison with the full-resolution and training-free acceleration baselines, we fine-tune on synthetic images generated by full-resolution Z-Image with 50 denoising steps using a set of 5K MS-COCO prompts that is disjoint from the 5K MS-COCO validation prompts used for evaluation. Training is conducted on our internal cluster of 8 NVIDIA H100 GPUs, each with 80~GB of VRAM, at an effective batch size of 8 and takes approximately 2 hours. We keep the optimizer and all other LoRA fine-tuning hyperparameters at the default settings of the \texttt{DiffSynth Studio} library, including a constant learning rate of 1e-4. Fine-tuning compute scales linearly with the number of sampled training images, training iterations, and batch size. In our experiments we fix the synthetic fine-tuning set to 5K prompts and train LoRA adapters for 2000 iterations, so increasing the fine-tuning dataset size would only affect compute through additional sampled training examples and does not change the inference-time cost of our method.

All latent-image generation experiments are conducted at a native resolution of $1024^2$ pixels, which is the default configuration for both FLUX.1-dev and Z-Image. The progressive resolution settings for the $S = 2$ and $S = 3$ cases are the same as in the FLUX.1-dev training-free acceleration experiments in Sec.~\ref{subsec:image_latent}.

Further model fine-tuning allows a more aggressive error tolerance $\delta$, leading to an even higher speedup. For fine-tuning experiments, we additionally evaluate $\delta=0.05$, which is larger than the default $\delta=0.01$. As in Sec.~\ref{subsec:image_latent}, we organize the comparison into three wall-clock speedup tiers, all calculated relative to the 50-step Z-Image baseline on a single H100 GPU (we adopt H100 wall-clock measurements for Z-Image and all subsequent pixel-space image generation and video generation experiments):

\begin{itemize}
    \item \textbf{Small speedup}: the 50-step native full-resolution Z-Image baseline, compared against our training-free acceleration and LoRA fine-tuned variants at $S=2$ and $S=3$ ($\delta=0.01$, 50 steps), as well as our LoRA fine-tuned variant with the configuration ($\delta=0.05$, $S=2$, 50 steps), yielding $1.65\times$ to $2.01\times$ speedup;
    \item \textbf{Medium speedup}: the 10-step native full-resolution Z-Image baseline, used as an iso-compute reference for our training-free acceleration and LoRA fine-tuned variants at $S=2$ and $S=3$ with reduced denoising steps, yielding approximately $5\times$ speedup; and
    \item \textbf{Large speedup}: the 7-step native full-resolution Z-Image baseline, used as an iso-compute reference for our training-free acceleration and LoRA fine-tuned variants at $S=2$ and $S=3$ with further reduced denoising steps, yielding approximately $7\times$ to $8\times$ speedup.
\end{itemize}
For quantitative evaluation, we use the same metric evaluation settings as Sec.~\ref{subsec:image_latent}, following the RALU evaluation protocol, and report the numbers in Table~\ref{tab:zimage-supp}. We show extensive qualitative comparisons in Figures~\ref{fig:zimage-1}--\ref{fig:zimage-3}. Across all three speedup tiers, our training-free acceleration variants outperform the reduced-step Z-Image baselines, and our spectral-transformation-based LoRA fine-tuning further closes the gap to the 50-step full-resolution baseline while attaining up to $7.81\times$ wall-clock speedup. In particular, our LoRA fine-tuned model at $S=3$ outperforms the training-free acceleration variant at $S=2$ while improving efficiency, and our LoRA fine-tuned model at $S=2$ either with reduced denoising steps or at a larger $\delta=0.05$ attains image quality comparable to the training-free acceleration $S=2$, 50-step setting, providing additional efficiency at no quality cost.


\begin{table*}[t!]
    \centering
    \caption{\textbf{Quantitative comparisons on latent-space image generation with fine-tuning.} We evaluate Z-Image at $1024^2$ resolution and group methods into wall-clock speedup tiers. Speedup is calculated relative to the 50-step Z-Image baseline wall-clock runtime on a single H100 GPU. \textcolor{blue}{\textbf{T}} and \textcolor{red}{\textbf{S}} denote temporal and spatial acceleration methods, respectively. All metrics are evaluated using the same protocol as Table~\ref{tab:flux-main}.}
    \label{tab:zimage-supp}
    \resizebox{\textwidth}{!}{
    \begin{scriptsize}
    \begin{tabular}{cccccccccc}
         \toprule
        \multirow{2}{*}{Method} & \multirow{2}{*}{Accel.} & \multirow{2}{*}{Latency (s) $\downarrow$} & \multirow{2}{*}{Speedup (s) $\uparrow$} & \multirow{2}{*}{TFLOPs $\downarrow$} & Overall & \multicolumn{2}{c}{Image quality} & \multicolumn{2}{c}{Text alignment} \\
        \cmidrule(lr){6-10}
        & & & & & ImageReward $\uparrow$ & CLIP-IQA $\uparrow$ & NIQE $\downarrow$ & T2I-Comp. $\uparrow$ & GenEval $\uparrow$  \\
        \midrule
        Z-Image (50 steps) & - & 21.25 & 1.00$\times$ & 4941.23 & 0.965 & 0.700 & 5.41 & 0.731 & 0.745  \\
        \cmidrule(lr){1-10}
        Z-Image (32 steps) & \textcolor{blue}{\textbf{T}} & 13.59 & 1.56$\times$ & 3166.62 & \underline{0.957} & 0.697 & \textbf{5.44} & 0.686 & 0.725  \\
        \textbf{Ours (TF, $\delta=0.01$, $S=2$)}     & \textcolor{red}{\textbf{S}}       & 12.90 & 1.65$\times$ & 3132.03 & 0.904 & 0.688 & 5.87 & 0.658 & \underline{0.730} \\
        \textbf{Ours (TF, $\delta=0.01$, $S=3$)}     & \textcolor{red}{\textbf{S}}       & \underline{12.19} & \underline{1.74$\times$} & \underline{2871.09} & 0.875 & 0.690 & \underline{5.59} & 0.650 & 0.682 \\
        \textbf{Ours (LoRA, $\delta=0.01$, $S=2$)}     & \textcolor{red}{\textbf{S}}       & 12.87 & 1.65$\times$ & 3132.03 & \textbf{0.982} & \textbf{0.699} & 5.72 & \textbf{0.725} & \textbf{0.731} \\
        \textbf{Ours (LoRA, $\delta=0.01$, $S=3$)}     & \textcolor{red}{\textbf{S}}       & 12.23 & 1.74$\times$ & \underline{2871.09} & 0.954 & \underline{0.697} & 5.75 & \underline{0.717} & 0.728 \\
        \textbf{Ours (LoRA, $\delta=0.05$, $S=2$)}     & \textcolor{red}{\textbf{S}}       & \textbf{10.59} & \textbf{2.01$\times$} & \textbf{2436.19} & 0.919 & 0.684 & 6.12 & 0.661 & 0.718 \\
        \cmidrule(lr){1-10}
        Z-Image (10 steps) & \textcolor{blue}{\textbf{T}} & 4.26 & 4.99$\times$ & 997.68 & 0.851 & 0.659 & 5.95 & \underline{0.678} & \underline{0.705} \\
        \textbf{Ours (TF, $\delta=0.01$, $S=2$)}     & \textcolor{red}{\textbf{S}}       & \underline{4.22} & \underline{5.04$\times$} & \underline{962.95} & \underline{0.860} & \underline{0.668} & 6.17 & 0.677 & 0.693 \\
        \textbf{Ours (TF, $\delta=0.01$, $S=3$)}     & \textcolor{red}{\textbf{S}}       & \textbf{4.01} & \textbf{5.30$\times$} & \textbf{875.97} & 0.827 & 0.662 & \underline{5.78} & 0.676 & 0.655 \\
        \textbf{Ours (LoRA, $\delta=0.01$, $S=2$)}     & \textcolor{red}{\textbf{S}}       &  4.24 & 5.01$\times$ & \underline{962.95} & \textbf{0.923} & \textbf{0.683} & \textbf{5.58} & \textbf{0.706} & \textbf{0.738} \\
        \cmidrule(lr){1-10}
        Z-Image (7 steps) & \textcolor{blue}{\textbf{T}} & 3.56 & 5.97$\times$ & 701.92 & 0.763 & 0.631 & \textbf{6.14} & 0.648 & 0.667 \\
        \textbf{Ours (TF, $\delta=0.01$, $S=2$)}     & \textcolor{red}{\textbf{S}}       & 2.75 & 7.73$\times$ & \underline{678.77} & \underline{0.804} & \underline{0.641} & 6.46 & \underline{0.665} & \underline{0.680} \\
        \textbf{Ours (TF, $\delta=0.01$, $S=3$)}     & \textcolor{red}{\textbf{S}}       & \textbf{2.63} & \textbf{8.08$\times$} & \textbf{609.18} & 0.759 & 0.630 & \underline{6.16} & 0.658 & 0.645 \\
        \textbf{Ours (LoRA, $\delta=0.01$, $S=2$)}     & \textcolor{red}{\textbf{S}}       & \underline{2.72} & \underline{7.81$\times$} & \underline{678.77} & \textbf{0.918} & \textbf{0.658} & 6.57 & \textbf{0.683} & \textbf{0.744} \\
        \bottomrule
    \end{tabular}
    \end{scriptsize}
    }
\end{table*}


% \subsection{Spectral Transform-based Fine-tuning for Latent- and Pixel-Space Image Generation}
% \label{subsec:image_pixel}

\paragraph{Additional pixel-space fine-tuning results.} For pixel-space image generation, we evaluate our method on PixelGen-XXL/16 T2I~\citep{pixelgen} at its native resolution of $512^2$ pixels. We use the same fine-tuning data construction and LoRA training procedure as the Z-Image latent-space experiments above, with synthetic fine-tuning images generated by full-resolution PixelGen with 25 denoising steps from a set of 5K MS-COCO prompts disjoint from the 5K MS-COCO validation prompts used for evaluation. Since PixelGen is pretrained with $x_0$ prediction rather than flow matching, we convert the spectral-transform velocity targets from the fine-tuning procedure in Sec.~\ref{sec:method:lora} to $x_0$ targets and supervise the LoRA adapter with the converted $x_0$.

We use $S=2$ with $\delta=0.01$ and, following Sec.~\ref{subsec:image_latent}, organize the comparison into two wall-clock speedup tiers, all calculated relative to the 25-step PixelGen baseline on a single H100 GPU:

\begin{itemize}
    \item \textbf{Small speedup}: the 25-step native full-resolution PixelGen baseline, used as the reference for the reduced 13-step PixelGen variant and our training-free acceleration and LoRA fine-tuned variants at $S=2$ (25 steps), yielding approximately $1.5\times$ to $1.9\times$ speedup;
    \item \textbf{Medium speedup}: the 8-step native full-resolution PixelGen variant, used as an iso-compute reference for our training-free acceleration and LoRA fine-tuned variants at $S=2$ with reduced denoising steps, yielding approximately $2\times$ to $3\times$ speedup.
\end{itemize}

For quantitative evaluation, we follow the same protocol used in latent-space image generation and report the numbers in Table~\ref{tab:pixelgen-main}. We show extensive qualitative comparisons in Figures~\ref{fig:pixelgen-1}--\ref{fig:pixelgen-3}. Because PixelGen does not natively support a high-fidelity $256^2$ resolution generation, our training-free acceleration variant falls short of the reduced-step PixelGen baselines on overall image quality across both speedup tiers. Our spectral-transformation-based LoRA fine-tuning closes this gap: at the same FLOPs as the training-free acceleration variant, the LoRA fine-tuned model matches or surpasses the reduced-step PixelGen baseline on most image-quality and prompt-alignment metrics, including ImageReward, CLIP-IQA, and T2I-CompBench. The qualitative comparisons in Figures~\ref{fig:pixelgen-1}--\ref{fig:pixelgen-3} corroborate this trend, showing that our LoRA fine-tuned model's outputs match the visual quality of the 25-step native-resolution PixelGen baseline at substantially higher throughput.


\begin{table*}[t!]
    \centering
    \caption{\textbf{Quantitative comparisons on pixel-space image generation.} The number in parentheses next to PixelGen indicates the total number of inference steps. $\uparrow / \downarrow$ denotes that a higher / lower metric is favourable. Speedup is calculated relative to the 25-step PixelGen baseline wall-clock runtime on a single H100 GPU. \textcolor{blue}{\textbf{T}} and \textcolor{red}{\textbf{S}} denote temporal and spatial acceleration, respectively.}
    \label{tab:pixelgen-main}
    \resizebox{\textwidth}{!}{
    \begin{scriptsize}
    \begin{tabular}{cccccccccc}
         \toprule
        \multirow{2}{*}{Method} & \multirow{2}{*}{Accel.} & \multirow{2}{*}{Latency (s) $\downarrow$} & \multirow{2}{*}{Speedup (s) $\uparrow$} & \multirow{2}{*}{TFLOPs $\downarrow$} & Overall & \multicolumn{2}{c}{Image quality} & \multicolumn{2}{c}{Text alignment} \\
        \cmidrule(lr){6-10}
        & & & & & ImageReward $\uparrow$ & CLIP-IQA $\uparrow$ & NIQE $\downarrow$ & T2I-Comp. $\uparrow$ & GenEval $\uparrow$  \\
        \midrule
        PixelGen (25 steps) & - & 0.48 & 1.00$\times$ & 65.36 & 0.921 & 0.734 & 5.95 & 0.574 & 0.794 \\
        \cmidrule(lr){1-10}
        PixelGen (13 steps) & \textcolor{blue}{\textbf{T}} & \textbf{0.25} & \textbf{1.92$\times$} & \underline{34.16} & \underline{0.886} & \underline{0.726} & \underline{5.89} & \underline{0.571} & \underline{0.781} \\
        \textbf{Ours (TF, $\delta=0.01$, $S = 2$)}     & \textcolor{red}{\textbf{S}}       & \underline{0.30} & \underline{1.60$\times$} & \textbf{33.72} & 0.799 & 0.718 & 6.10 & 0.568 & \textbf{0.782} \\
        \textbf{Ours (LoRA, $\delta=0.01$, $S = 2$)}   & \textcolor{red}{\textbf{S}}       & 0.31 & 1.55$\times$ & \textbf{33.72} & \textbf{0.913} & \textbf{0.728} & \textbf{5.87} & \textbf{0.580} & 0.776 \\
        \cmidrule(lr){1-10}
        PixelGen (8 steps) & \textcolor{blue}{\textbf{T}} & \textbf{0.16} & \textbf{3.00$\times$} & \underline{21.16} & \underline{0.858} & \underline{0.715} & \textbf{5.79} & \underline{0.572} & 0.756 \\
        \textbf{Ours (TF, $\delta=0.01$, $S = 2$)}     & \textcolor{red}{\textbf{S}}       & \underline{0.21} & \underline{2.29$\times$} & \textbf{20.75} & 0.779 & 0.713 & 6.02 & 0.567 & \underline{0.765} \\
        \textbf{Ours (LoRA, $\delta=0.01$, $S = 2$)}   & \textcolor{red}{\textbf{S}}       & 0.22 & 2.18$\times$ & \textbf{20.75} & \textbf{0.908} & \textbf{0.724} & \underline{5.85} & \textbf{0.577} & \textbf{0.770} \\
        \bottomrule
    \end{tabular}
    \end{scriptsize}
    }
\end{table*}


\clearpage
\begin{figure}
    \centering
    \includegraphics[width=\linewidth]{supp_figures/flux_supp_1.pdf}
    \caption{\textbf{Qualitative comparisons on latent-space image generation.} We compare our method against default-step generation, reduced-step native-resolution generation on FLUX.1-dev \citep{black2024flux}, and RALU~\citep{ralu}, a state-of-the-art acceleration baseline matched to similar speedups. Our method outperforms both baselines. FLUX.1-dev with reduced steps degrades image quality and exhibits over-saturation artifacts, while RALU introduces noticeable noise artifacts.}
    \label{fig:flux-1}
\end{figure}

\clearpage
\begin{figure}
    \centering
    \includegraphics[width=\linewidth]{supp_figures/flux_supp_2.pdf}
    \caption{\textbf{Qualitative comparisons on latent-space image generation.} We compare our method against default-step generation, reduced-step native-resolution generation on FLUX.1-dev \citep{black2024flux}, and RALU~\citep{ralu}, a state-of-the-art acceleration baseline matched to similar speedups. Our method outperforms both baselines. FLUX.1-dev with reduced steps degrades image quality and exhibits over-saturation artifacts, while RALU introduces noticeable noise artifacts.}
    \label{fig:flux-2}
\end{figure}

\clearpage
\begin{figure}
    \centering
    \includegraphics[width=\linewidth]{supp_figures/flux_supp_3.pdf}
    \caption{\textbf{Qualitative comparisons on latent-space image generation.} We compare our method against default-step generation, reduced-step native-resolution generation on FLUX.1-dev \citep{black2024flux}, and RALU~\citep{ralu}, a state-of-the-art acceleration baseline matched to similar speedups. Our method outperforms both baselines. FLUX.1-dev with reduced steps degrades image quality and exhibits over-saturation artifacts, while RALU introduces noticeable noise artifacts.}
    \label{fig:flux-3}
\end{figure}

\iffalse
\clearpage
\begin{figure}
    \centering
    \includegraphics[width=\linewidth]{supp_figures/flux_supp_4.pdf}
    \caption{\textbf{Qualitative comparisons on latent-space image generation.} We compare our method against default-step generation, reduced-step native-resolution generation on FLUX.1-dev \citep{black2024flux}, and RALU~\citep{ralu}, a state-of-the-art acceleration baseline matched to similar speedups. Our method outperforms both baselines. FLUX.1-dev with reduced steps degrades image quality and exhibits over-saturation artifacts, while RALU introduces noticeable noise artifacts.}
    \label{fig:flux-4}
\end{figure}

\clearpage
\begin{figure}
    \centering
    \includegraphics[width=\linewidth]{supp_figures/flux_supp_5.pdf}
    \caption{\textbf{Qualitative comparisons on latent-space image generation.} We compare our method against default-step generation, reduced-step native-resolution generation on FLUX.1-dev \citep{black2024flux}, and RALU~\citep{ralu}, a state-of-the-art acceleration baseline matched to similar speedups. Our method outperforms both baselines. FLUX.1-dev with reduced steps degrades image quality and exhibits over-saturation artifacts, while RALU introduces noticeable noise artifacts.}
    \label{fig:flux-5}
\end{figure}
\fi

\clearpage
\begin{figure}
     \centering
     \includegraphics[width=\linewidth]{supp_figures/zimage_supp_1.pdf}
     \caption{\textbf{Qualitative comparisons on latent-space image generation (fine-tuned).} We compare our method against default-step generation, reduced-step native-resolution generation on Z-Image \citep{zimage} matched to similar speedups. Our fine-tuned model (Ours$^*$) achieves even higher image quality compared to our training-free acceleration variant and outperforms the reduced-step baseline.}
     \label{fig:zimage-1}
\end{figure}

\clearpage
\begin{figure}
     \centering
     \includegraphics[width=\linewidth]{supp_figures/zimage_supp_2.pdf}
     \caption{\textbf{Qualitative comparisons on latent-space image generation (fine-tuned).} We compare our method against default-step generation, reduced-step native-resolution generation on Z-Image \citep{zimage} matched to similar speedups. Our fine-tuned model (Ours$^*$) achieves even higher image quality compared to our training-free acceleration variant and outperforms the reduced-step baseline.}
     \label{fig:zimage-2}
\end{figure}

\clearpage
\begin{figure}
     \centering
     \includegraphics[width=\linewidth]{supp_figures/zimage_supp_3.pdf}
     \caption{\textbf{Qualitative comparisons on latent-space image generation (fine-tuned).} We compare our method against default-step generation, reduced-step native-resolution generation on Z-Image \citep{zimage} matched to similar speedups. Our fine-tuned model (Ours$^*$) achieves even higher image quality compared to our training-free acceleration variant and outperforms the reduced-step baseline.}
     \label{fig:zimage-3}
\end{figure}


\iffalse
\clearpage
\begin{figure}
     \centering
     \includegraphics[width=\linewidth]{supp_figures/zimage_supp_4.pdf}
     \caption{\textbf{Qualitative comparisons on latent-space image generation (fine-tuned).} We compare our method against default-step generation, reduced-step native-resolution generation on Z-Image \citep{zimage} matched to similar speedups. Our fine-tuned model (Ours$^*$) achieves even higher image quality compared to our training-free acceleration variant and outperforms the reduced-step baseline.}
     \label{fig:zimage-4}
\end{figure}

\clearpage
\begin{figure}
     \centering
     \includegraphics[width=\linewidth]{supp_figures/zimage_supp_5.pdf}
     \caption{\textbf{Qualitative comparisons on latent-space image generation (fine-tuned).} We compare our method against default-step generation, reduced-step native-resolution generation on Z-Image \citep{zimage} matched to similar speedups. Our fine-tuned model (Ours$^*$) achieves even higher image quality compared to our training-free acceleration variant and outperforms the reduced-step baseline.}
     \label{fig:zimage-5}
\end{figure}
\fi

\clearpage
\begin{figure}
    \centering
    \includegraphics[width=0.67\linewidth]{supp_figures/pixelgen_supp_1.pdf}
    \caption{\textbf{Qualitative comparisons on pixel-space image generation.} We compare our method against default-step generation and reduced-step native-resolution generation on PixelGen~\citep{pixelgen}, matched to comparable speedups. An asterisk (Ours$^*$) marks the fine-tuned model. Our method achieves similar quality to full-resolution generation while attaining higher speedups.}
    \label{fig:pixelgen-1}
\end{figure}

\clearpage
\begin{figure}
    \centering
    \includegraphics[width=0.67\linewidth]{supp_figures/pixelgen_supp_2.pdf}
    \caption{\textbf{Qualitative comparisons on pixel-space image generation.} We compare our method against default-step generation and reduced-step native-resolution generation on PixelGen~\citep{pixelgen}, matched to comparable speedups. An asterisk (Ours$^*$) marks the fine-tuned model. Our method achieves similar quality to full-resolution generation while attaining higher speedups.}
    \label{fig:pixelgen-2}
\end{figure}

\clearpage
\begin{figure}
    \centering
    \includegraphics[width=0.67\linewidth]{supp_figures/pixelgen_supp_3.pdf}
    \caption{\textbf{Qualitative comparisons on pixel-space image generation.} We compare our method against default-step generation and reduced-step native-resolution generation on PixelGen~\citep{pixelgen}, matched to comparable speedups. An asterisk (Ours$^*$) marks the fine-tuned model. Our method achieves similar quality to full-resolution generation while attaining higher speedups.}
    \label{fig:pixelgen-3}
\end{figure}

\iffalse
\clearpage
\begin{figure}
    \centering
    \includegraphics[width=0.67\linewidth]{supp_figures/pixelgen_supp_4.pdf}
    \caption{\textbf{Qualitative comparisons on pixel-space image generation.} We compare our method against default-step generation and reduced-step native-resolution generation on PixelGen~\citep{pixelgen}, matched to comparable speedups. An asterisk (Ours$^*$) marks the fine-tuned model. Our method achieves similar quality to reduced-step generation while attaining higher speedups.}
    \label{fig:pixelgen-4}
\end{figure}


\clearpage
\begin{figure}
    \centering
    \includegraphics[width=0.67\linewidth]{supp_figures/pixelgen_supp_5.pdf}
    \caption{\textbf{Qualitative comparisons on pixel-space image generation.} We compare our method against default-step generation and reduced-step native-resolution generation on PixelGen~\citep{pixelgen}, matched to comparable speedups. An asterisk (Ours$^*$) marks the fine-tuned model. Our method achieves similar quality to reduced-step generation while attaining higher speedups.}
    \label{fig:pixelgen-5}
\end{figure}
\fi
